\begin{tikzpicture}[scale=1.0, font=\small]
  % Shared body of one panel; the channel drawing differs per panel and is added after.
  \newcommand{\mospanel}[3]{%
    \fill[ecteal, opacity=0.10] (0,0) rectangle (5.0,1.9);
    \draw[ecgray, line width=0.8pt] (0,0) rectangle (5.0,1.9);
    \node[note, anchor=south west, font=\footnotesize] at (0.12,0.12){$p$ 衬底};
    \fill[ecred, opacity=0.20] (0.4,1.3) rectangle (1.5,1.9);
    \fill[ecred, opacity=0.20] (3.5,1.3) rectangle (4.6,1.9);
    \draw[ecgray, line width=0.7pt] (0.4,1.3) rectangle (1.5,1.9);
    \draw[ecgray, line width=0.7pt] (3.5,1.3) rectangle (4.6,1.9);
    \node[font=\footnotesize] at (0.95,1.6){$n^{+}$};
    \node[font=\footnotesize] at (4.05,1.6){$n^{+}$};
    \fill[eclight, opacity=0.55] (1.5,1.9) rectangle (3.5,2.12);
    \fill[ecgray, opacity=0.45] (1.5,2.12) rectangle (3.5,2.5);
    \draw[ecgray, line width=0.9pt] (0.95,1.9) -- (0.95,3.0) node[above]{S};
    \draw[ecgray, line width=0.9pt] (2.5,2.5) -- (2.5,3.0) node[above]{G};
    \draw[ecgray, line width=0.9pt] (4.05,1.9) -- (4.05,3.0) node[above]{D};
    \node[note, anchor=north, font=\footnotesize] at (2.5,4.05){#2};
    \node[anchor=south, font=\small] at (2.5,4.35){#1};
    \node[note, anchor=north, align=center, font=\footnotesize] at (2.5,-0.3){#3};
  }

  % ---------- (a) cut-off ----------
  \begin{scope}
    \mospanel{截止区}{$V_{GS}<V_{TH}$}{栅下只有耗尽层，\\ S 与 D 之间没有导电通路}
    \draw[eclight, dashed, line width=0.9pt] (1.5,1.6) -- (3.5,1.6);
    \node[note, anchor=north, font=\footnotesize] at (2.5,1.55){耗尽层};
  \end{scope}

  % ---------- (b) triode ----------
  \begin{scope}[xshift=6.4cm]
    \mospanel{三极管区}{$V_{GS}>V_{TH}$，$V_{DS}$ 很小}{沟道厚度均匀，\\ 像一个受 $V_{GS}$ 控制的电阻}
    \fill[ecblue, opacity=0.6] (1.5,1.74) rectangle (3.5,1.9);
    \node[ecblue, font=\footnotesize, anchor=north] at (2.5,1.68){反型沟道};
  \end{scope}

  % ---------- (c) saturation ----------
  \begin{scope}[xshift=12.8cm]
    \mospanel{饱和区}{$V_{DS}>V_{GS}-V_{TH}$}{漏端沟道夹断，\\ $I_D$ 几乎不再随 $V_{DS}$ 变化}
    \fill[ecblue, opacity=0.6]
      (1.5,1.9) -- (3.35,1.9) .. controls (2.7,1.87) .. (1.5,1.74) -- cycle;
    \draw[ecred, line width=0.9pt] (3.35,1.9) -- (3.5,1.9);
    \node[ecred, font=\footnotesize, anchor=north east] at (3.75,1.66){夹断点};
    \draw[helper, ecred] (3.4,1.88) -- (3.5,1.7);
  \end{scope}

  \node[note, anchor=north, align=center] at (8.9,-1.35)
    {MOS 是纯电压控制器件：栅极隔着氧化层「静电地」控制沟道，直流栅电流为零 ——
     这正是它比 BJT 更适合做大规模集成的根本原因};
\end{tikzpicture}
